% 修订跟踪命令集合
% 使用方法：
% \added[id=作者,remark=说明]{新增内容}
% \deleted[id=作者,remark=说明]{删除内容}
% \replaced[id=作者,remark=说明]{新内容}{旧内容}

\usepackage{xcolor}
\usepackage{soul}

% 定义颜色
\definecolor{addedcolor}{rgb}{0,0.6,0}
\definecolor{deletedcolor}{rgb}{0.8,0,0}
\definecolor{replacedcolor}{rgb}{0,0,0.8}

% 新增内容命令
\newcommand{\added}[2][]{%
  \textcolor{addedcolor}{\ulcorner{#2}\urcorner}%
  \ifx&#1\else%
    \textcolor{gray}{\tiny[#1]}%
  \fi%
}

% 删除内容命令
\newcommand{\deleted}[2][]{%
  \textcolor{deletedcolor}{\sout{#2}}%
  \ifx&#1\else%
    \textcolor{gray}{\tiny[#1]}%
  \fi%
}

% 替换内容命令
\newcommand{\replaced}[3][]{%
  \deleted[#1]{#3} $\rightarrow$ \added[#1]{#2}%
}

% 批准修订命令（将修订设为正式）
\newcommand{\accept}[1]{#1}

% 拒绝修订命令（移除修订标记）
\newcommand{\reject}[1]{}

% 简化的使用示例
% \added{这是新增内容}
% \deleted[reviewer]{这是删除内容}
% \replaced[editor]{新内容}{旧内容}